\begin{chaptersummary}
  \item A schema has no version number, so a field you regret is a field you
        live with. Every decision in this chapter is about making the next
        change additive rather than breaking.
  \item A global node id is the type name and the key, base64 encoded:
        \code{U2Vzc2lvbjox} is \code{Session:1}. One \code{[NodeResolver]}
        method per type rewrites its \code{id} to \code{ID!}, makes it
        implement \code{Node}, and gives every object in the graph one way to
        be refetched. The encoding is not a secret, and a well-formed id of
        the wrong type resolves to the wrong object with no error at all.
  \item \code{@deprecated} is advisory and that is its whole value: the field
        keeps answering, and introspection stops offering it, so tooling
        migrates clients while nothing breaks. The reason string is the only
        message the person removing the field will read, so it names the
        replacement. It may not go on a required argument, and the schema
        refuses to build if it does.
  \item A failure the caller should act on belongs in the payload, not in the
        top-level \code{errors} array where it is a string nobody can branch
        on. Two \code{[Error<T>]} attributes generate a union of error types
        and an \code{Error} interface they both implement: the union tells a
        client every failure that exists today, and the interface lets
        tomorrow's arrive without breaking that client. The answer is 200,
        and the double-booking in this chapter wrote nothing.
  \item A non-null field is a promise the specification collects on by
        propagating the null upward until it reaches a position that may hold
        one. One session pointing at a missing speaker emptied a list of four,
        at status 200, with one error. Promise a field only when you own the
        value and something enforces it.
\end{chaptersummary}
